.bp odd
.pn 1
.cs 5
.mt 5
.mb 6
.pl 66
.ll 65
.po 10
.hm 2
.fm 2
.fi
.he
.ft                                1-%
.ce 2
.sh
Section 1
.qs
.sp
.sh
Introduction
.qs
.sp 2
.pp 0
Digital Research PL/I is an implementation of PL/I for microcomputers that use
the 8080, 8086, 8088, or similar processor.  It is formally based on American
National Standard X3.74 PL/I General Purpose Subset (Subset G).  Subset G has
the formal structure of the full PL/I language, but in some ways it is a new
language and in many ways an improved language compared to full PL/I.
.pp
PL/I Subset G is easy to learn and use.  It is a highly portable language
because its design usually ensures hardware independence.  It is also more
efficient and cost effective.  Programs written in PL/I Subset G are easier to
implement, document, and maintain.
.sp 2
.sh
1.1  Documentation Set
.qs
.he PL/I Reference Manual                      1.1  Documentation Set
.tc    1.1  Documentation Set
.pp
The \c 
.ul
PL/I Language Reference Manual
.qu
\ presents a detailed but concise description of the PL/I programming 
language.  It is not a tutorial on how to program in PL/I; rather, it is a 
functional description of the language, its syntax, and semantics.  
This manual is a reference document that supplements Digital 
Research's \c 
.ul
PL/I Language Programmer's Guide.
.qu
.ix PL/I keywords
.ix PL/I syntax
.ix data attributes
.ix variables
.ix identifiers
.pp
The \c 
.ul
PL/I Language Programmer's Guide
.qu
\ includes sample programs that illustrate many of the features of 
PL/I, as well as the mechanical aspects of compiling and linking 
programs.  If you have not programmed in PL/I before, read 
the Programmer's Guide first, while cross-referencing specific topics 
in the Reference Manual.  If you are already an experienced PL/I 
programmer, you might want to read the Reference Manual only.
.pp
The \c 
.ul
PL/I Language Command Summary
.qu
\ lists all the PL/I keywords and statement forms, data attributes, and 
error messages.  It also contains a summary of the commands for the 
compiler.
.he PL/I Reference Manual                               1.2  Notation
.bp
.sh
1.2  Notation
.qs
.tc    1.2  Notation
.ix special characters
.ix control characters
.ix BIF, acronymn for built-in function
.ix built-in functions
.pp
The following notational conventions appear throughout this document:
.sp 2
.in 5
.ti -2
o Words in capital letters are PL/I keywords.
.sp
.ti -2
o Words in lower-case letters or in a combination of lower-case letters 
and digits separated by a hyphen represent variable information for you 
to select.  These words are described or defined more explicitly in the 
text.
.sp
.ti -2 
o Example statements are given in lower-case.  
.sp
.ti -2
o The vertical bar | indicates alternatives.
.sp
.ti -2
o b represents a blank character.
.sp
.ti -2
o Square brackets [] enclose options.
.sp
.ti -2
o Ellipses (...) indicate that the immediately preceding item can occur 
once, or any number of times in succession.
.sp
.ti -2
o Except for the special characters listed above, all other 
punctuation and special characters represent the actual occurrence of 
those characters.
.sp
.ti -2
o Within the text, the symbol CTRL represents a control character.  
Thus, CTRL-C means control-C.  In a PL/I source program listing or any 
listing that shows example console interaction, the symbol ^ represents 
a control character.
.sp
.ti -2
o The acronym BIF means built-in function.
.sp
.ti -2
o Everything that you type at the keyboard and that appears 
on the screen is in colored type.
.sp 2
.ce
End of Section 1
.he
.nx twoa.tex


 
